\newpage
\section{Aplikacja webowa}
\label{sec:aplikacja}

Polecenie \kod{moderator web} (w Dockerze uruchamiane automatycznie) startuje aplikację Gradio pod adresem
\url{http://localhost:7860}. Nagłówek pokazuje używane modele i to, czy Ollama odpowiada, więc brakujący model widać
od razu. Każda zakładka jest osobnym plikiem w \kod{src/moderator/web/tabs/} i korzysta z tej samej klasy
\kod{Moderator} co terminal. Wynik oceny ma postać karty: kolor i plakietka decyzji, pasek stopnia naruszenia,
uzasadnienie, kategorie i -- w trybie RAG -- cytowany paragraf razem z jego treścią. Przyciski nad nagłówkiem
przełączają język interfejsu (PL / EN) oraz motyw jasny i ciemny; wybór zapamiętuje przeglądarka. Zmienia się tylko
interfejs, a komentarze, odpowiedzi modelu i regulamin zostają bez zmian.

\zrzut{web_01_single_comment}{Zakładka \emph{Single comment}: groźba oceniona w trybie RAG z cytowanym paragrafem.}

\zrzut{web_02_image}{Zakładka \emph{Image}: tekst odczytany ze zrzutu ekranu i jego ocena.}

\zrzut{web_03_batch}{Zakładka \emph{Batch}: statystyki, wykres, tabela i pliki wynikowe dla 40 komentarzy.}

\zrzut{web_04_chat}{Zakładka \emph{Chat about a file}: pytanie o groźby, dopytanie z historią i zapisane pliki.}

\zrzut{web_05_rules}{Zakładka \emph{Rules (RAG)}: wyszukiwanie paragrafów i cały regulamin.}

\zrzut{web_06_evaluation}{Zakładka \emph{Evaluation}: wyniki \kod{moderator evaluate}.}
